\documentclass[10pt]{article}
\usepackage[a4paper,margin=24mm]{geometry}
\usepackage{amsmath,amssymb,amsthm,mathtools}
\usepackage{booktabs,array,microtype}
\usepackage[hidelinks]{hyperref}
\providecommand{\CRevisionRound}{2}
\ifdefined\pdfvariable
  \pdfvariable suppressoptionalinfo 767
\fi
\newtheorem{theorem}{Theorem}
\newtheorem{proposition}[theorem]{Proposition}
\newtheorem{remark}[theorem]{Remark}
\newcommand{\Part}{\mathfrak P}
\newcommand{\rank}{\operatorname{rk}}
\newcommand{\fall}[2]{(#1)_{#2}}
\title{Parallel Binary Refinement of Labelled Set Partitions:\
Exact Kernels, Spectrum, and the Last-Collision Threshold}
\author{Route-A Dynamics Working Note C301}
\date{2 September 2026}

\begin{document}
\maketitle

\begin{abstract}
We solve a synchronous fragmentation chain on the labelled set partitions of
$[n]$.  At every step each label receives a fresh fair bit and each current
block is refined by its two nonempty bit fibres.  A coupling by binary words
gives the full $t$-step kernel from every starting partition, the complete
law of the number of blocks, and the exact absorption-time distribution.
The transition matrix has eigenvalues $2^{k-n}$ with Stirling multiplicities
$S(n,k)$.
\ifnum\CRevisionRound>0
A refinement flag and a genuine lowering argument give a squarefree
annihilating polynomial, closing the diagonalizability gap that triangularity
alone would leave.  The absorption window is the birthday window:
$n^2/2^t\to\lambda$ gives the limit $e^{-\lambda/2}$, with its dyadic lattice
boundary kept explicit.
\fi
\ifnum\CRevisionRound>1
An exact, independently checked finite certificate accompanies the analytic
proof.  A hostile Route-A audit returns
$(\mathrm{A0},\ldots,\mathrm{A4})=(\mathrm{FAIL},\ldots,\mathrm{FAIL})$;
the finite Markov determinant is not promoted to arithmetic data.
\fi
\end{abstract}

\section{The labelled parallel-fragmentation chain}

Let $n\geq1$ and let $\Part_n$ be the set of labelled set partitions of
$[n]=\{1,\ldots,n\}$.  Write $\pi\preceq\sigma$ when $\sigma$ refines
$\pi$.  Starting from $\Pi_0=\pi$, every label independently receives a
fresh bit in $\{0,1\}$.  In each block of $\pi$, labels with equal bits stay
together; empty fibres are deleted.  The resulting partition is $\Pi_1$.
Updates use fresh independent bits.

This is a synchronous chain: all blocks are marked in parallel.  Labels are
part of the state.  In particular, quotienting by block sizes produces a
different chain and different spectral multiplicities.  Put
\[
  \fall qr=q(q-1)\cdots(q-r+1),\qquad \fall q0=1,
\]
and let $S(n,k)$ denote a Stirling number of the second kind.  For
$\pi\preceq\sigma$, let $r_B(\pi,\sigma)$ be the number of $\sigma$-blocks
contained in $B\in\pi$.

\section{Complete law}

\begin{theorem}[all-parameter solution]\label{thm:main}
For every $n\geq1$, $t\geq0$, and $\pi,\sigma\in\Part_n$, with $q=2^t$,
\begin{equation}\label{eq:t-kernel}
 K_n^t(\pi,\sigma)=
 \mathbf 1_{\{\pi\preceq\sigma\}}
 \prod_{B\in\pi}\frac{\fall q{r_B(\pi,\sigma)}}{q^{|B|}}.
\end{equation}
Consequently, the one-step probability equals $2^{|\pi|-n}$ precisely when
$\sigma$ refines $\pi$ and each $\pi$-block contains at most two
$\sigma$-blocks, and equals zero otherwise.

From the one-block initial partition, for every $\sigma\in\Part_n$ and
$1\leq k\leq n$,
\begin{align}
 \Pr\{\Pi_t=\sigma\}&=\frac{\fall {2^t}{|\sigma|}}{2^{tn}},
 \label{eq:partition-law}\\
 \Pr\{|\Pi_t|=k\}&=S(n,k)\frac{\fall {2^t}k}{2^{tn}},
 \label{eq:block-law}\\
 \mathbb E|\Pi_t|&=2^t\bigl[1-(1-2^{-t})^n\bigr].
 \label{eq:mean-blocks}
\end{align}
Let $T_n=\min\{t:\Pi_t\text{ is discrete}\}$.  Then
\begin{align}
 \Pr\{T_n\leq t\}&=\frac{\fall {2^t}n}{2^{tn}},
 \label{eq:cdf}\\
 \Pr\{T_n=t\}&=\frac{\fall {2^t}n}{2^{tn}}
 -\mathbf 1_{\{t\geq1\}}\frac{\fall {2^{t-1}}n}{2^{(t-1)n}},
 \label{eq:mass}\\
 \mathbb ET_n&=\sum_{t\geq0}
 \left[1-\frac{\fall {2^t}n}{2^{tn}}\right].
 \label{eq:mean-time}
\end{align}
The series converges.  The characteristic polynomial, source spectral
determinant, and power traces are
\begin{align}
 \chi_{K_n}(x)&=\prod_{k=1}^n
    (x-2^{k-n})^{S(n,k)},\label{eq:charpoly}\\
 \det(I-zK_n)&=\prod_{k=1}^n
    (1-z2^{k-n})^{S(n,k)},\label{eq:det}\\
 \operatorname{tr}(K_n^t)&=\sum_{k=1}^nS(n,k)2^{t(k-n)}.
 \label{eq:trace}
\end{align}
Moreover,
\begin{equation}\label{eq:annihilator}
 \prod_{k=1}^n(K_n-2^{k-n}I)=0,
\end{equation}
so $K_n$ is diagonalizable over $\mathbb Q$.

Finally, if $n_j\to\infty$, $t_j$ are integers, and
$n_j^2/2^{t_j}\to\lambda\in(0,\infty)$, then
\begin{equation}\label{eq:critical}
 \Pr\{T_{n_j}\leq t_j\}\longrightarrow e^{-\lambda/2}.
\end{equation}
Thus $T_n$ has a tight $O(1)$ window around $2\log_2n$.
\end{theorem}

\section{Binary-word proof}

After $t$ updates, attach to label $i$ the word formed by its successive
bits.  These words are independent and uniform in a set of size $q=2^t$.
Inside each starting block $B$, two labels are together at time $t$ exactly
when their words agree.  To realize $r_B$ prescribed terminal fibres one
must injectively assign words to those fibres.  There are $\fall q{r_B}$
assignments.  Starting blocks use disjoint label sets, hence their counts
multiply, proving \eqref{eq:t-kernel}.  Taking $t=1$ gives the stated
one-step kernel.

For a one-block start, a fixed $k$-block partition has probability
$\fall qk/q^n$.  There are $S(n,k)$ such labelled partitions, yielding
\eqref{eq:partition-law} and \eqref{eq:block-law}; the identity
$\sum_k S(n,k)\fall qk=q^n$ also verifies normalization.  Alternatively,
$|\Pi_t|$ is the number of occupied boxes after $n$ independent throws into
$q$ boxes.  Summing the $q$ occupancy indicators proves
\eqref{eq:mean-blocks}.

Absorption means that all $n$ words are distinct, which proves
\eqref{eq:cdf}; differencing successive distribution functions gives
\eqref{eq:mass}.  The tail-sum formula gives \eqref{eq:mean-time}.  A union
bound gives
\[
  \Pr\{T_n>t\}\leq \binom n2 2^{-t},
\]
so the tail series converges.  For $n=1$, all formulas reduce to $T_1=0$.
The contract begins at $n=1$; the empty-set convention $n=0$ is not used.

\ifnum\CRevisionRound>0
\section{The spectral flag and its missing step}

Order states by their number of blocks.  Refinement makes $K_n$ triangular.
A $k$-block state stays unchanged exactly when every block receives a
constant bit; the $2^k$ successful assignments among $2^n$ have probability
\(\lambda_k=2^{k-n}\).  Hence the diagonal rank-$k$ block is
$\lambda_k I_{S(n,k)}$, proving \eqref{eq:charpoly};
\eqref{eq:det} and \eqref{eq:trace} follow once diagonalizability is closed.

Triangularity alone is not that closure.  Let $V_k$ be the span of basis
vectors indexed by partitions of rank at most $k$, with the conventional
column action of the displayed transition matrix.  The spaces
$0\subset V_1\subset\cdots\subset V_n$ are invariant, and
$(K_n-\lambda_k I)V_k\subseteq V_{k-1}$.  Since all polynomial factors in
$K_n$ commute, applying the factors in descending rank lowers this flag one
step at a time.  Their product kills $V_n$, proving
\eqref{eq:annihilator}.  The roots $\lambda_1,\ldots,\lambda_n$ are distinct,
so the minimal polynomial is squarefree and $K_n$ is diagonalizable over
$\mathbb Q$.  This also justifies the trace formula without any hidden
Jordan-block assumption.

\section{Birthday threshold and lattice boundary}

Write $q_j=2^{t_j}$.  Under the hypotheses of \eqref{eq:critical},
$n_j/q_j\to0$ and
\begin{align*}
 \log\frac{\fall {q_j}{n_j}}{q_j^{n_j}}
 &=\sum_{r=0}^{n_j-1}\log(1-r/q_j)\\
 &=-\frac{n_j(n_j-1)}{2q_j}
   +O\!\left(\frac{n_j^3}{q_j^2}\right)
 \longrightarrow-\lambda/2.
\end{align*}
Exponentiation proves \eqref{eq:critical}.  The same product, together with
the union bound above, yields tightness around $2\log_2n$.

There is an essential lattice qualification.  Because $t$ is integer and
$q=2^t$, the quantity $n^2/2^{\lfloor2\log_2n+c\rfloor}$ need not converge
without controlling its dyadic phase.  We therefore assert the limit only
along sequences for which that ratio converges; no phase-free continuous law
for $T_n-2\log_2n$ is claimed.
\fi

\ifnum\CRevisionRound>1
\section{Exact certificate and hostile boundary audit}

The machine-readable certificate is not used to replace the proof.  It
enumerates all $278$ labelled partitions for $1\leq n\leq6$, reconstructs
all $1{,}860$ nonzero transition cells independently, and checks $81$ exact
time-law rows for $1\leq n\leq9$ and $0\leq t\leq8$.  A separate SymPy
implementation verifies characteristic polynomials, determinants,
squarefree annihilators, eigenspace dimensions, and $t$-step kernels through
$n=5$.  Deterministic replay and a mutation suite protect the serialization,
formula, scope, and evaluation contracts.  These finite checks are regression
evidence only; Theorem~\ref{thm:main} is global because of the preceding
word coupling and flag proof.

The closest source owner is the Hopf-power/rock-breaking framework of
Diaconis, Pang, and Ram~\cite{DPR}.  Their unlabelled integer-partition
rock-breaking chain and labelled-graph examples make clear that binary
breaking, absorption, and Hopf diagonalization are established mechanisms.
This note claims neither priority nor a new general Hopf-algebra theorem.  Its
closed contribution is a self-contained labelled-set-partition formulation
with the complete kernel, occupancy law, exact last-collision distribution,
elementary spectral-flag proof, and executable audit in one contract.

The Route-A tuple is
\[
 (A0,A1,A2,A3,A4)=(\mathrm{FAIL},\mathrm{FAIL},\mathrm{FAIL},
 \mathrm{FAIL},\mathrm{FAIL}),
 \qquad \texttt{ROUTE\_A\_REJECTED}.
\]
There is no arithmetic local datum or Euler factor (A0).  Refinement has no
nonconstant directed cycle---only self-loops---and therefore no primitive
orbit repetition law (A1).  Ordinary integer time and the scale $2^t$ are not
an arithmetic logarithmic clock (A2).  Although \eqref{eq:det} is exact, it
is a finite source Markov polynomial, not a target completed determinant or
functional equation (A3).  Finally, rational diagonalizability is not a
self-adjoint target-zero lift (A4).  Route B is locked under the literal
\texttt{NO\_BAD\_EULER\_OR\_ROOT\_NUMBER} scope.

Three model mutations are outside the theorem: assigning one shared bit per
block prevents splitting; biased label bits replace the uniform injection
count; and quotienting by block sizes changes the state space and Stirling
multiplicities.  These are different systems, not harmless reparametrizations.

\paragraph{AI-use statement.}
An AI assistant helped draft prose and executable checks.  The formulas were
rederived independently, finite claims were regenerated from exact arithmetic,
and the final logical and scope assertions remain the responsibility of the
release author.
\fi

\begin{thebibliography}{1}
\bibitem{DPR}
P. Diaconis, C. Y. A. Pang, and A. Ram,
``Hopf algebras and Markov chains: two examples and a theory,''
\emph{Journal of Algebraic Combinatorics} (2014),
doi:10.1007/s10801-013-0456-7; arXiv:1206.3620.
\end{thebibliography}

\end{document}
